\begin{theorem}[Wild quadratic case]\label{mod:cases-wildquadratic-main}
Let \(K/F\) be a prime-cyclic extension of nonarchimedean local fields of
degree \(2\), with residue characteristic \(2\), residue degree one, and
positive lower break \(t\).  Fix an integral uniformizer \(\pi_K\) which
generates \(\mathcal O_K\) over \(\mathcal O_F\).  Let
\(\Delta_F,\Delta_K\) be local-constant functions satisfying the proved
finite-sum specification \texttt{IsDeltaFiniteLocalConstant}.

For arbitrary exact data \(\chi_F\) and \(\psi_F\), let
\(\chi_{\min}\) be the least-conductor member of the complete
norm-character orbit of \(\chi_F\).  Higher stationary input is required
only under the implication
\[
  1<m_F(\chi_{\min}).
\]
In that event, \texttt{WildQuadraticCommonHigherData} retains exactly the
q-free information shared by every branch: the coefficient view indexed at
the actual conductor \(m_F(\chi_{\min})\), its identification with
\(\chi_{\min}\) and \(\psi_F\), the
\texttt{FirstMainComputationalData}, \texttt{FirstMainPhaseData}, real
\texttt{ExactQuadraticPhaseAssembly}, and common coordinate compatibility.
The branch constructors then add either exact actual-even rows or the
refined phase-row compatibility.

The higher package is branch-sensitive:
\[
\begin{array}{c|l}
\texttt{allEven}&
 \text{q-free all-even coefficient view, four actual even rows, and}
 \\
 &\texttt{ExactQuadraticPhaseAssembly.CompleteCorrectionData};\\[1mm]
\texttt{refined}&
 \text{an actual \texttt{PositivePolarSource}, a normalized refinement tied}
 \\
 &\text{to that source, the existing refined coefficient view and row}
 \\
 &\text{compatibility, refinement assembly, and mixed-boundary continuation.}
\end{array}
\]
Thus the \texttt{allEven} constructor has no refinement field, while the
\texttt{refined} constructor cannot be formed without provenance from one
of the seven actual positive-polar rows.  The all-even/positive-polar
disjointness proof makes both impossible combinations unrepresentable.
Only when a refined row is the exceptional odd boundary and \(2\ne0\) in
\(F\) is the package required to supply the boundary residue data, the two
actual selected refinement-correction records, their boundary compatibility,
and their selection equalities.  No boundary witness is requested in equal
characteristic two or in a nonboundary row.

Then Lean proves the literal First Main identity
\[
 \Delta_K(\chi_F\circ N_{K/F},\psi_F\circ\operatorname{Tr}_{K/F})
 \prod_{\mu\in S(K/F)}\Delta_F(\mu,\psi_F)
 =
 \prod_{\mu\in S(K/F)}\Delta_F(\mu\chi_F,\psi_F),
\]
where both products range over the complete norm-character group, including
the identity.

The proof first uses whole-orbit minimality and dispatches on the actual
conductor of \(\chi_{\min}\).  Conductor zero invokes
\texttt{firstMain\_wildEndpoint\_zero}, and conductor one invokes
\texttt{firstMain\_wildEndpoint\_one}; the finite-Delta hypotheses identify
their concrete \texttt{deltaFinite} products with the displayed
\(\Delta_F,\Delta_K\) products.  No stationary object is introduced in
either endpoint branch.

For \(m_F(\chi_{\min})>1\), Lean first eliminates the branch-sensitive
higher package.  The \texttt{allEven} constructor is sent directly to
\texttt{wildQuadratic\_nonboundary\_allEven}, so no \(\mathfrak q\) is
manufactured for the three constant rows.  In the \texttt{refined}
constructor, the coefficient row is either the exceptional
\texttt{boundaryOdd} row or its complement.  The complement lies among the
nine disjoint nonboundary rows---equivalently
\[
 m<T,\qquad m=T\text{ with }T\text{ even},\qquad\text{or}\qquad T<m,
 \quad T=t+1,
\]
and is discharged by \texttt{wildQuadratic\_nonboundary} with the derived
whole-orbit minimality proof.  In the exceptional row, the coefficient
classifier gives \(T=t+1\) odd.  If \(\operatorname{char}F=2\), the
equal-characteristic break theorem gives \(t\) odd, while oddness of
\(t+1\) gives that \(t\) is not odd, a contradiction.  Thus break parity is
used in this direction and no equal-characteristic boundary data are
assumed.  If \(2\ne0\) in \(F\), the exact error formula and precisely the
supplied boundary records are passed to
\texttt{wildQuadratic\_boundary}.  These alternatives are exhaustive and
pairwise disjoint.

Finally, the proved equivalence between error one and the fixed-character
First Main identity converts the higher cancellation result.  Membership of
\(\chi_{\min}\) in the original norm-character orbit and the already-proved
reindexing invariance of both complete products transport the identity back
to \(\chi_F\).
\LeanFile{LanglandsFirstMainLemma/Cases/WildQuadratic/Main.lean}
\lean{LanglandsFirstMainLemma.WildQuadraticCommonHigherData}
\lean{LanglandsFirstMainLemma.WildQuadraticHigherData}
\lean{LanglandsFirstMainLemma.WildQuadraticHigherData.allEven}
\lean{LanglandsFirstMainLemma.WildQuadraticHigherData.refined}
\lean{LanglandsFirstMainLemma.firstMain_wildQuadratic}
\leanok
\SourceText{Proposition prop:wild-quadratic-cancellation and the final Wild quadratic case corollary.}
\uses{mod:cases-wildquadratic-nonboundary,mod:cases-wildquadratic-boundary,mod:cases-wildendpoints,mod:parameters-minimalorbitstationary}
\FormalizationContract
The theorem does not claim to reconstruct the higher stationary assembly
from an arbitrary continuous character.  It consumes that exact completed
assembly only after proving that the chosen minimal representative has
conductor greater than one.  In particular, the endpoint branches require
no \texttt{WildQuadraticHigherData}.

The public conclusion is \texttt{FirstMainIdentity}, not merely an isolated
finite-field phase equality.  The higher package fixes its coefficient
index to the actual multiplicative conductor of the selected minimal
representative and retains the \(t+1\) norm-character boundary.  Its
\texttt{allEven} branch is q-free; its \texttt{refined} branch contains an
actual positive-polar source and a refinement normalized from that source,
and asks for mixed-characteristic boundary compatibility only in the genuine
\texttt{boundaryOdd} branch.  All stationary representatives, correction
records, and residual phases are exactly those supplied by the completed
nonboundary and boundary interfaces; none is reconstructed or re-assumed.
\end{theorem}
